\begin{table}[!ht]
\centering
\caption{\label{tab:rf_xgboost_metrics}Random Forest and XGBoost: Out-of-Sample Performance by Horizon}
\footnotesize
\setlength{\tabcolsep}{3.5pt}

\begin{tabular}{llrrrrrrrrr}
\toprule
Model & Horizon & RMSE & MAE & R2 & DA & Naive DA & Train RMSE & Train MAE & Train R2 & Train DA\\
\midrule
Random Forest & 1d & 0.01660 & 0.01065 & -0.03685 & 0.34095 & 0.27619 & 0.00903 & 0.00550 & 0.69418 & 0.60486\\
Random Forest & 3d & 0.02698 & 0.01978 & -0.05024 & 0.50095 & 0.32000 & 0.01537 & 0.01118 & 0.65356 & 0.82984\\
Random Forest & 5d & 0.03333 & 0.02579 & -0.06204 & 0.47238 & 0.30286 & 0.01817 & 0.01317 & 0.69659 & 0.87750\\
XGBoost & 1d & 0.01632 & 0.00974 & -0.00288 & 0.28571 & 0.27619 & 0.01633 & 0.00954 & 0.00000 & 0.31888\\
XGBoost & 3d & 0.02653 & 0.01923 & -0.01508 & 0.52190 & 0.32000 & 0.02349 & 0.01677 & 0.19113 & 0.66635\\
XGBoost & 5d & 0.03271 & 0.02481 & -0.02306 & 0.44190 & 0.30286 & 0.03183 & 0.02304 & 0.06882 & 0.50763\\
\bottomrule
\multicolumn{11}{p{\linewidth}}{\rule{0pt}{1em}\textit{Note.} DA = Directional Accuracy (fraction of test-set observations where the predicted sign matches the realised sign). Naive DA uses the previous day's return (lag1) as the prediction. R2 is out-of-sample and may be negative if the model performs worse than the mean.}\\
\end{tabular}
\end{table}
